\documentclass[12pt, a4paper]{book}

% ─── Core packages ────────────────────────────────────────────────────────────
\usepackage[T1]{fontenc}
\usepackage[utf8]{inputenc}
\usepackage[charter, cal=cmcal]{mathdesign}
\usepackage[scaled=0.92]{helvet}
\usepackage{geometry}
\geometry{
  a4paper,
  top=30mm, bottom=30mm,
  inner=32mm, outer=25mm,
  marginparwidth=18mm
}

% ─── Mathematics ──────────────────────────────────────────────────────────────
\usepackage{amsmath, amsthm}

% ─── Chemistry ────────────────────────────────────────────────────────────────
\usepackage[version=4]{mhchem}

% ─── Graphics & colour ────────────────────────────────────────────────────────
\usepackage{graphicx}
\usepackage[dvipsnames,svgnames,table]{xcolor}
\usepackage{tikz}
\usetikzlibrary{arrows.meta, shapes.geometric, positioning,
                calc, backgrounds, decorations.pathmorphing,
                decorations.markings, patterns}
\usepackage{pgfplots}
\pgfplotsset{compat=1.18}
\usepackage{tcolorbox}
\tcbuselibrary{skins, breakable, theorems}

% ─── Layout ───────────────────────────────────────────────────────────────────
\usepackage{titlesec}
\usepackage{fancyhdr}
\usepackage{enumitem}
\usepackage{booktabs}
\usepackage{array}
\usepackage{multicol}
\usepackage{marginnote}
\usepackage{mdframed}
\usepackage{wrapfig}
\usepackage{caption}
\usepackage{subcaption}
\usepackage{parskip}
\usepackage{siunitx}
\usepackage[colorlinks=true, linkcolor=BrickRed,
            urlcolor=NavyBlue, citecolor=ForestGreen]{hyperref}

% ══════════════════════════════════════════════════════════════════════════════
% COLOUR PALETTE  (identical to Module 1.1)
% ══════════════════════════════════════════════════════════════════════════════
\definecolor{JourneyColor}{HTML}{065F46}   % forest green — journey headings
\definecolor{ModuleColor}{HTML}{0F766E}    % deep teal — module headings
\definecolor{PuzzleColor}{HTML}{B45309}    % deep amber — puzzle boxes
\definecolor{DiscoveryColor}{HTML}{D97706} % golden amber — discovery boxes
\definecolor{WarningColor}{HTML}{E11D48}   % coral red — "model breaks" boxes
\definecolor{ProblemColor}{HTML}{0F766E}   % dark teal — problem sets
\definecolor{ReflectColor}{HTML}{047857}   % green — reflection boxes
\definecolor{HistoryColor}{HTML}{9A3412}   % brown — history notes
\definecolor{LightBg}{HTML}{F8FAFC}        % near-white background

% ══════════════════════════════════════════════════════════════════════════════
% TCOLORBOX ENVIRONMENTS
% ══════════════════════════════════════════════════════════════════════════════
\newtcolorbox{puzzlebox}[1][]{
  enhanced, breakable,
  colback=PuzzleColor!6, colframe=PuzzleColor,
  fonttitle=\bfseries\sffamily, title={The Puzzle},
  attach boxed title to top left={yshift=-2mm, xshift=6mm},
  boxed title style={colback=PuzzleColor, colframe=PuzzleColor},
  borderline west={3pt}{0pt}{PuzzleColor}, #1}

\newtcolorbox{discoverybox}[1][]{
  enhanced, breakable,
  colback=DiscoveryColor!8, colframe=DiscoveryColor,
  fonttitle=\bfseries\sffamily, title={Discovery},
  attach boxed title to top left={yshift=-2mm, xshift=6mm},
  boxed title style={colback=DiscoveryColor, colframe=DiscoveryColor},
  borderline west={3pt}{0pt}{DiscoveryColor}, #1}

\newtcolorbox{problembox}[1][]{
  enhanced, breakable,
  colback=ProblemColor!5, colframe=ProblemColor,
  fonttitle=\bfseries\sffamily, title={Problems},
  attach boxed title to top left={yshift=-2mm, xshift=6mm},
  boxed title style={colback=ProblemColor, colframe=ProblemColor},
  borderline west={3pt}{0pt}{ProblemColor}, #1}

\newtcolorbox{deepbox}[1][]{
  enhanced, breakable,
  colback=JourneyColor!5, colframe=JourneyColor,
  fonttitle=\bfseries\sffamily, title={Deep Problems},
  attach boxed title to top left={yshift=-2mm, xshift=6mm},
  boxed title style={colback=JourneyColor, colframe=JourneyColor},
  borderline west={4pt}{0pt}{JourneyColor}, #1}

\newtcolorbox{historybox}[1][]{
  enhanced, breakable,
  colback=HistoryColor!7, colframe=HistoryColor!60,
  fonttitle=\bfseries\sffamily\small, title={Historical Note},
  attach boxed title to top left={yshift=-2mm, xshift=6mm},
  boxed title style={colback=HistoryColor!80, colframe=HistoryColor}, #1}

\newtcolorbox{reflectionbox}[1][]{
  enhanced, breakable,
  colback=ReflectColor!6, colframe=ReflectColor,
  fonttitle=\bfseries\sffamily, title={Reflection},
  attach boxed title to top left={yshift=-2mm, xshift=6mm},
  boxed title style={colback=ReflectColor, colframe=ReflectColor},
  borderline west={3pt}{0pt}{ReflectColor}, #1}

\newtcolorbox{modelbreaks}[1][]{
  enhanced, breakable,
  colback=WarningColor!6, colframe=WarningColor,
  fonttitle=\bfseries\sffamily, title={Where This Model Breaks},
  attach boxed title to top left={yshift=-2mm, xshift=6mm},
  boxed title style={colback=WarningColor, colframe=WarningColor},
  borderline west={4pt}{0pt}{WarningColor}, #1}

\newtcolorbox{namedconcept}[2][]{
  enhanced,
  colback=ModuleColor!6, colframe=ModuleColor,
  fonttitle=\bfseries\sffamily\large, title={#2},
  attach boxed title to top left={yshift=-2mm, xshift=6mm},
  boxed title style={colback=ModuleColor, colframe=ModuleColor}, #1}

% ══════════════════════════════════════════════════════════════════════════════
% THEOREM ENVIRONMENTS
% ══════════════════════════════════════════════════════════════════════════════
\theoremstyle{definition}
\newtheorem{observation}{Observation}[chapter]
\newtheorem{question}{Question}[chapter]
\newtheorem{experiment}{Experiment}[chapter]

% ══════════════════════════════════════════════════════════════════════════════
% SECTION FORMATTING
% ══════════════════════════════════════════════════════════════════════════════
\titleformat{\chapter}[display]
  {\normalfont\sffamily\huge\bfseries\color{JourneyColor}}
  {\textcolor{JourneyColor!60}{\chaptertitlename\ \thechapter}}
  {10pt}{\Huge}[\vspace{2pt}\textcolor{JourneyColor}{\titlerule[1.5pt]}]

\titleformat{\section}
  {\normalfont\sffamily\Large\bfseries\color{ModuleColor}}
  {\thesection}{1em}{}[\textcolor{ModuleColor!40}{\titlerule[0.6pt]}]

\titleformat{\subsection}
  {\normalfont\sffamily\large\bfseries\color{JourneyColor!80}}
  {\thesubsection}{1em}{}

\titleformat{\subsubsection}
  {\normalfont\sffamily\normalsize\bfseries}
  {\thesubsubsection}{1em}{}

% ══════════════════════════════════════════════════════════════════════════════
% HEADER / FOOTER
% ══════════════════════════════════════════════════════════════════════════════
\pagestyle{fancy}
\fancyhf{}
\fancyhead[LE]{\small\sffamily\textcolor{JourneyColor}{\leftmark}}
\fancyhead[RO]{\small\sffamily\textcolor{ModuleColor}{\rightmark}}
\fancyfoot[C]{\small\thepage}
\renewcommand{\headrulewidth}{0.4pt}

% ══════════════════════════════════════════════════════════════════════════════
% CUSTOM COMMANDS
% ══════════════════════════════════════════════════════════════════════════════
\newcommand{\chem}[1]{\ce{#1}}
\newcommand{\important}[1]{\textcolor{ModuleColor}{\textbf{#1}}}
\newcommand{\newterm}[1]{\textit{\textcolor{DiscoveryColor}{#1}}}
\newcommand{\sidenote}[1]{%
  \marginnote{\footnotesize\sffamily\textcolor{JourneyColor!70}{#1}}}

% ══════════════════════════════════════════════════════════════════════════════
% DOCUMENT
% ══════════════════════════════════════════════════════════════════════════════
\begin{document}
\setcounter{chapter}{0}   % Journey 1 is chapter 1

% We skip the title page and TOC here — those live in the master file.
% This file compiles standalone for review.

\chapter[Journey 1 — What Is Stuff?]
        {Journey One\\[8pt]
         \textcolor{ModuleColor}{\large What Is Stuff?\ \ Matter and Its Patterns}}

% Fast-forward to Module 1.2
\setcounter{section}{1}   % so next \section is 1.2

% ══════════════════════════════════════════════════════════════════════════════
%  MODULE 1.2 — BURNING, DISSOLVING, REACTING: THE FIRST CLUES
% ══════════════════════════════════════════════════════════════════════════════
\section{Module 1.2 — Burning, Dissolving, Reacting: The First Clues}
\label{sec:mod1.2}

\subsection*{Motivation}

In Module~1.1 we learned to classify matter:
elements, compounds, mixtures.
We noticed that when elements combine into a compound,
something \emph{new} appears — a substance whose properties
are entirely unlike its components.

But we have not yet asked the hardest question:

\begin{center}
\begin{tcolorbox}[
  enhanced, colback=LightBg, colframe=JourneyColor!40,
  width=0.86\textwidth, arc=3pt,
  fontupper=\sffamily
]
\centering
When a chemical reaction happens,\\[4pt]
\textbf{what is actually going on?}\\[6pt]
\small
When wood burns, where does the wood \emph{go}?\\
When iron rusts, what is iron combining with?\\
When limestone dissolves in acid, what is released?\\[4pt]
\normalsize
For two thousand years, these questions had no satisfying answer.\\
Then, in the span of a single generation,\\
careful measurement changed everything.
\end{tcolorbox}
\end{center}

\vspace{0.5em}

This module follows the intellectual detective story
that led from confused guesses about fire and rust
to one of the most powerful laws in all of science:
\textbf{the conservation of mass}.

% ─── Section: The Phlogiston Theory ──────────────────────────────────────────
\subsection{The Theory That Ruled for a Century}

Before we reach the correct answer, we must spend time with
the \emph{wrong} one — because understanding why a wrong theory
is compelling teaches us more about science than jumping
directly to the truth.

\subsubsection{The Phlogiston Hypothesis}

In the early 18th century, the German chemist Georg Ernst Stahl
proposed a theory that seemed to explain combustion, rusting,
and many other reactions in a single elegant framework.

His idea: every combustible substance contains a hidden substance
called \newterm{phlogiston}.
When something burns, it releases phlogiston into the air.
When something rusts (which Stahl saw as a slow form of burning),
it also releases phlogiston slowly.
When a metal ore is heated with charcoal, the charcoal
donates its phlogiston to the ore, restoring the metal.

\begin{center}
\begin{tikzpicture}[
  every node/.style={font=\sffamily\small},
  arrow/.style={-{Stealth[length=5pt]}, thick, color=HistoryColor!80}
]
  % Combustion row
  \node[draw=HistoryColor!60, fill=HistoryColor!8, rounded corners=3pt,
        minimum width=3.0cm, minimum height=0.9cm, align=center]
        (wood) at (0,0) {Wood\\(rich in phlogiston)};
  \draw[arrow] (wood) -- ++(3.2,0)
    node[above, midway, font=\sffamily\footnotesize]{burns};
  \node[draw=HistoryColor!60, fill=HistoryColor!8, rounded corners=3pt,
        minimum width=3.0cm, minimum height=0.9cm, align=center]
        (ash) at (4.5,0) {Ash\\(phlogiston-free)};
  \node[draw=JourneyColor!50, fill=JourneyColor!8, rounded corners=3pt,
        minimum width=2.4cm, minimum height=0.7cm, align=center,
        dashed]
        (phi) at (2.4,-1.4) {phlogiston\\released to air};
  \draw[arrow, dashed, color=JourneyColor!60] (1.6,-0.4) -- (phi.north west);

  % Rusting row
  \node[draw=HistoryColor!60, fill=HistoryColor!8, rounded corners=3pt,
        minimum width=3.0cm, minimum height=0.9cm, align=center]
        (iron) at (0,-3.0) {Iron\\(contains phlogiston)};
  \draw[arrow] (iron) -- ++(3.2,0)
    node[above, midway, font=\sffamily\footnotesize]{rusts slowly};
  \node[draw=HistoryColor!60, fill=HistoryColor!8, rounded corners=3pt,
        minimum width=3.0cm, minimum height=0.9cm, align=center]
        (rust) at (4.5,-3.0) {Iron oxide (rust)\\(phlogiston-free)};
  \node[font=\sffamily\footnotesize\itshape, text=JourneyColor!70]
        at (2.4,-4.0) {phlogiston slowly lost to air};
\end{tikzpicture}
\end{center}

\vspace{0.5em}

The theory was internally consistent.
It explained why things needed air to burn
(air absorbs phlogiston, and air saturated with phlogiston
cannot absorb more — hence a candle dies in a closed jar).
It explained smelting of metals from ores.
It gave chemists a shared vocabulary.

It was taught in every university in Europe for nearly a century.

\begin{historybox}
\textbf{Why Smart People Believed Phlogiston (1700s)}

Georg Stahl's theory was not foolish — it was the best available
explanation given the tools of the time.
No one had yet invented a reliable way to collect and weigh gases.
Reactions happened in open vessels where gases escaped unnoticed.

The critical failure of phlogiston theory was that it treated
\emph{fire} as the release of a substance,
when fire is actually the \emph{combination} with a substance (oxygen).
In a sense, Stahl had the direction of the arrow backwards.

This type of error — getting the mechanism exactly reversed —
appears repeatedly in the history of science.
It is only detectable by precise, quantitative measurement.
\end{historybox}

% ─── Puzzle: the weight anomaly ──────────────────────────────────────────────
\subsection{The Crack in the Theory: Mass Increases on Burning}

The first serious blow to phlogiston theory came from a simple
but deeply uncomfortable observation.

\begin{puzzlebox}
When you burn a piece of wood in an open container,
the ash weighs \emph{less} than the wood.
This fits phlogiston theory: the phlogiston was released,
so the ash is lighter.

But when you heat a piece of tin or lead in air,
the resulting oxide (called a \emph{calx}) weighs
\emph{more} than the original metal.

\begin{center}
\begin{tikzpicture}[every node/.style={font=\sffamily\small}]
  % Left: wood burning
  \node[draw=JourneyColor!50, fill=JourneyColor!8,
        rounded corners=3pt, minimum width=3cm,
        minimum height=0.8cm, align=center]
        (w1) at (0,0) {Wood: \SI{100}{\gram}};
  \draw[-{Stealth}, thick, color=DiscoveryColor]
        (0,-0.5) -- (0,-1.5) node[midway, right]{burns in air};
  \node[draw=JourneyColor!50, fill=JourneyColor!8,
        rounded corners=3pt, minimum width=3cm,
        minimum height=0.8cm, align=center]
        (w2) at (0,-2) {Ash: \SI{5}{\gram} \quad (\emph{lighter})};

  % Right: tin heating
  \node[draw=WarningColor!60, fill=WarningColor!8,
        rounded corners=3pt, minimum width=3cm,
        minimum height=0.8cm, align=center]
        (t1) at (7,0) {Tin: \SI{100}{\gram}};
  \draw[-{Stealth}, thick, color=WarningColor]
        (7,-0.5) -- (7,-1.5) node[midway, right]{heated in air};
  \node[draw=WarningColor!60, fill=WarningColor!8,
        rounded corners=3pt, minimum width=3cm,
        minimum height=0.8cm, align=center]
        (t2) at (7,-2) {Tin oxide: \SI{127}{\gram} \quad (\emph{heavier!})};

  \node[font=\sffamily\footnotesize\itshape, text=JourneyColor!70]
        at (0,-3) {as expected by phlogiston};
  \node[font=\sffamily\footnotesize\bfseries\itshape, text=WarningColor]
        at (7,-3) {NOT explained by phlogiston};
\end{tikzpicture}
\end{center}

\begin{enumerate}[label=(\alph*), leftmargin=*]
  \item If phlogiston is released during burning,
        why would the product \emph{gain} mass?

  \item One attempted fix: ``phlogiston has negative mass.''
        A substance releases phlogiston, so the residue becomes heavier.
        This is logically consistent. Is it scientifically satisfying?
        What would you need to show to test this ``negative mass'' idea?

  \item What does the mass increase of tin oxide suggest
        about what might be \emph{added} to tin during heating?
\end{enumerate}
\end{puzzlebox}

\vspace{0.5em}

The phlogiston supporters invented increasingly desperate patches:
negative-mass phlogiston,
phlogiston absorbed from the air rather than released,
and so on.
Each patch made the theory more complex without explaining more.
This is a classical sign that a theory is failing.

% ─── Section: Lavoisier ──────────────────────────────────────────────────────
\subsection{Lavoisier: The Power of Weighing Everything}

Antoine-Laurent Lavoisier was not the first to notice
the mass-increase anomaly.
What made him extraordinary was his \emph{method}:
he refused to accept any explanation that could not be
quantitatively verified by weighing.

His key insight was deceptively simple:

\begin{center}
\begin{tcolorbox}[
  enhanced, colback=DiscoveryColor!8, colframe=DiscoveryColor,
  width=0.82\textwidth, arc=3pt,
  fontupper=\sffamily\large\bfseries, halign=center
]
If you weigh \emph{everything} that goes into a reaction\\
and \emph{everything} that comes out,\\[4pt]
\normalfont\normalsize\sffamily
including the gases, including the air,\\
including the vapour and the smoke,\\[4pt]
\bfseries
the total mass does not change.
\end{tcolorbox}
\end{center}

\vspace{0.5em}

This sounds straightforward.
In practice, it required Lavoisier to invent new apparatus:
sealed vessels to capture gases,
precision balances accurate to milligrams,
and careful accounting of every component.

\begin{historybox}
\textbf{Lavoisier and the Burning of Phosphorus (1772)}

One of Lavoisier's key early experiments:
he burned phosphorus in a sealed glass vessel sitting in water.
After burning, he found:
\begin{itemize}[nosep]
  \item The vessel was lighter than before
        (the phosphorus had combined with some of the air).
  \item But when the seal was broken,
        air rushed \emph{in} — exactly enough to restore the original mass.
  \item The product weighed more than the phosphorus.
\end{itemize}

Conclusion: phosphorus had combined with part of the air.
The ``missing'' mass was not missing at all — it was in the air
that had been consumed.

\begin{center}
\ce{4P + 5O2 -> P4O10}
\end{center}

This single experiment made it clear that combustion is \emph{combination}
with something from the air — not release of something from the fuel.
\end{historybox}

% ─── Section: Oxygen ─────────────────────────────────────────────────────────
\subsection{The Discovery of Oxygen}

If burning requires combining with \emph{part} of the air,
what is that part?

Around 1774, Carl Wilhelm Scheele in Sweden and
Joseph Priestley in England independently isolated a gas
by heating mercury(II) oxide:

\begin{center}
\ce{2HgO_{(s)} ->[\Delta] 2Hg_{(l)} + O2_{(g)}}
\end{center}

Priestley found the gas made a candle burn brighter
and made mice in a closed vessel survive longer.
He called it \emph{dephlogisticated air} — air with no phlogiston,
therefore able to absorb more from burning things.
He was still thinking inside the phlogiston framework.

Lavoisier repeated Priestley's experiment, measured everything
carefully, and reached a different conclusion.
He named the gas \emph{oxygène} — ``acid-former'' —
and showed it was a distinct element
that combined with substances during burning.

\begin{center}
\begin{tikzpicture}[
  every node/.style={font=\sffamily\small},
  arrow/.style={-{Stealth[length=5pt]}, thick}
]
  % Reactants
  \node[draw=JourneyColor, fill=JourneyColor!8, rounded corners=3pt,
        minimum width=2.2cm, minimum height=1cm, align=center]
        (fuel) at (0,0) {Fuel\\(\emph{e.g.} carbon)};
  \node[draw=ModuleColor, fill=ModuleColor!8, rounded corners=3pt,
        minimum width=2.2cm, minimum height=1cm, align=center]
        (oxy) at (3.5,0) {Oxygen\\from air};

  % Combine arrow
  \node at (1.75,0) {\Large $+$};
  \draw[arrow, color=DiscoveryColor] (4.7,0) -- (6.2,0)
    node[above, midway, font=\sffamily\footnotesize]{react};

  % Product
  \node[draw=DiscoveryColor!80, fill=DiscoveryColor!8, rounded corners=3pt,
        minimum width=3cm, minimum height=1cm, align=center]
        (prod) at (8,0) {Carbon dioxide\\+ heat + light};

  % Mass arrows
  \draw[arrow, color=HistoryColor, dashed]
    (0,-0.8) -- (0,-1.8) node[below, font=\sffamily\footnotesize]{mass: $m_1$};
  \draw[arrow, color=HistoryColor, dashed]
    (3.5,-0.8) -- (3.5,-1.8) node[below, font=\sffamily\footnotesize]{mass: $m_2$};
  \draw[arrow, color=HistoryColor, dashed]
    (8,-0.8) -- (8,-1.8) node[below, font=\sffamily\footnotesize]{mass: $m_1 + m_2$};
\end{tikzpicture}
\end{center}

The key distinction between Priestley's view and Lavoisier's view
is not about the facts — both observed the same gas.
It is about the \emph{framework} used to interpret the facts.
Priestley interpreted the new gas inside phlogiston theory.
Lavoisier built a new theory to fit the quantitative measurements.

\begin{observation}
\textbf{The lesson of Lavoisier vs Priestley:}
Two scientists can observe the \emph{same phenomenon}
and reach \emph{opposite conclusions} depending on
which framework they use to interpret it.

A better framework is not one that fits the current data better.
It is one that makes \emph{more precise, testable predictions}.
\end{observation}

% ─── Section: Conservation of Mass ───────────────────────────────────────────
\subsection{The Law of Conservation of Mass}

From hundreds of carefully weighed experiments,
Lavoisier arrived at the most fundamental quantitative law
in chemistry.

\end{document}